\documentclass{beamer}
\usepackage{amsmath}
\usepackage{amsfonts}
\usepackage{amssymb}
\usepackage{yhmath}
\usepackage{amsthm}
\usepackage{mathrsfs}
\usepackage{mathtools}
\usepackage{mathalfa}
\usepackage{upgreek}
\usepackage{yfonts}
\usetheme{Warsaw}
%\usetheme{Madrid}
%\usetheme{Bergen}
%\usetheme{AnnArbor}
%\usetheme{CambridgeUS}\setbeamercovered{dynamic}
%\usetheme{Hannover}
%\usetheme{Antibes}
%\usetheme{Copenhagen}
%\usetheme{Marburg} 
\usefonttheme{serif}
\begin{document}
\title{Adding Highly Generic Subsets of $\omega_2$}

\author{Rouholah Hoseini Naveh}
%{
%\setbeamercolor{background canvas}{bg=yellow}
\begin{frame}
\date{}
\maketitle

\begin{center}
Joint work with:\\
Dr Esfandiar Eslami\\
Dr Mohammad Golshani\\
\vspace{10mm}
1401/12/3
\end{center}
\end{frame}
\begin{frame}
\frametitle{CONSISTENCY AND INDEPENDENCE}
It all started on that day in December 1873 when Georg Cantor established that \textit{\textbf{{\large the continuum is not countable}}}

$$|X| < |\mathcal{P}(X)| \qquad \qquad \aleph_0, \aleph_1, \aleph_2, \dots$$
\pause
$$c= 2^{\aleph_0} = \aleph_?$$
$$2^{\aleph_{\alpha}}= \aleph_{\alpha+1}$$
\pause
\textbf{G\"{o}del 1939:} \ $Con(ZF) \ implies \ Con(ZFC + GCH)$
\vspace{5mm}

\textbf{Cohen 1963:} $ Con(ZF) \ implies \ Con(ZF + \neg AC)$

\hspace{25.5mm}$Con(ZFC) \ implies \ Con(ZFC + \neg CH)$
$$V[G] \models ZFC + 2^{\aleph_0}\geq \aleph_2$$
\end{frame}
\begin{frame}
\frametitle{FORCING}
\begin{block}{Forcing Notion}
$\mathbb{P}= \langle P, \leq_P \rangle \in V$ a poset of \textbf{conditions} with largest element

\hspace{20mm}$Add(\aleph_0, \aleph_2)= \{p \vdots (\omega \times \omega_2) \longrightarrow 2 \ | \ |p| < \aleph_0\}$

$p_1 \leq p_2$ \quad \textbf{{\large or}}\quad $p_1$ extends $p_2$ \quad \textbf{{\large or}}\quad $p_1$ is stronger than $p_2$

\textbf{{\large or}}\quad $p_1$ has more information than $p_2$ \ \textbf{{\large iff}} \ $p_2 \subseteq p_1$
\end{block}
\vspace{-2mm}
\begin{block}{Dense Open subsets}
$D \subseteq \mathbb{P}$ is dense open if
\begin{itemize}
\item $\forall p \in \mathbb{P} \exists d \in D (d \leq p)$
\item $d_1 \in D \land d_2 \leq d_1 \Rightarrow d_2 \in D$
\end{itemize}
\end{block}
\vspace{-2mm}
\begin{block}{Generic Set}
$G$ is Generic if
\begin{itemize}
	\item $G \subseteq \mathbb{P}$ is a filter
	\item $G$ meets every dense open subset of $\mathbb{P}$ that lies in $V$
\end{itemize}
\end{block}
\end{frame}
\begin{frame}
\frametitle{PRESERVING CARDINALS AND THE GCH}
Cohen showed that $(2^{\aleph_0})^{V[G]} \geq \aleph_2^{V}$

Does $\aleph_2^{V} = \aleph_2^{V[G]}$?
\pause
\begin{block}{Chain condition}
$\mathbb{P}$ is $\kappa.c.c.$ if every maximal antichain $A \subseteq \mathbb{P}$ has size $<\kappa$

\textbf{Cohen: If $\mathbb{P}$ is $\kappa.c.c.$ then forcing with $\mathbb{P}$ preserves cardinals $\geq \kappa$.}
\end{block}
\end{frame}
\begin{frame}
\frametitle{$\aleph_1$-PRESERVING FORCINGS}
\textbf{Recall: }Let $X$ be a set and $\kappa \leq |X|$ a cardinal
$$[X]^{\kappa}=\{Y \subseteq X \ | \ |Y|= \kappa\}$$
\begin{block}{Proper Forcings}
A forcing notion $\mathbb{P}$ is proper if for every infinite $X$ and every stationary set $S \subseteq [X]^{\leq \aleph_0}$, $S$ remains stationary in $V[G]$
\begin{itemize}
\item \textbf{Shelah: If $\mathbb{P}$ is proper then forcing with $\mathbb{P}$ preserves $\aleph_1$}
\end{itemize}
\end{block}
\end{frame}
\begin{frame}
\frametitle{ELEMENTARY SUBSTRUCTURES}
Let $\theta$ be a regular uncountable cardinal
\begin{itemize}
\item $H(\theta)=\{x: \ |TC(x)| < \theta \}$\\
\hspace{9mm}$=\{x: \ x \subseteq y \ \exists y(y \ \text{is transitive} \land |y| < \theta)\}$
\item $\langle H(\theta), \in \rangle \models ZFC-P$
\item If $\theta$ is inaccessible cardinal then $\langle H(\theta), \in \rangle \models ZFC$
\end{itemize}
\begin{block}{Elementary substructures}
$\mathfrak{A} \prec \mathfrak{B}$ iff $\mathfrak{A} \subseteq \mathfrak{B}$ and for all formulas $\varphi[x_1, \dots, x_n]$ of $\mathscr{L}$ and all $a_1, \dots, a_n \in A$, we have
$$\mathfrak{A}\models \varphi[a_1, \dots, a_n] \iff \mathfrak{B}\models \varphi[a_1, \dots, a_n]$$
\end{block}
Let $\mathfrak{B}$ be a model of power $\alpha$, let $|\mathscr{L}| \leq \beta \leq \alpha$, let $X \subseteq B$ and $|X| \leq \beta$ Then $\mathfrak{B}$ has an elementary submodel of power $\beta$ which contains $X$
\end{frame}
\begin{frame}
\frametitle{FORCING WITH ELEMENTARY SUBSTRUCTURES}
$$\mathcal{S}= \{M \in [H(\theta)]^{\aleph_0} : \langle M, \in, <_w \rangle \prec \langle H(\theta), \in, <_w \rangle \}$$
\vspace{-2mm}
\begin{block}{THE $\in$-COLLAPSE FORCING}
$\mathbb{P}_{\in}(\theta)$ is the set of all finite $\in$-chains of countable elementary submodels of $\langle H(\theta), \in, \leq_w \rangle $ with the inverse inclusion as the order
\end{block}
\begin{block}{MATRIX $\in$-COLLAPSE FORCING}
$\mathbb{P}_{\in}^{\mathcal{M}}= \{p \subset \mathcal{S} \ | \ |p| < \aleph_0 \}$ such that
\begin{itemize}
\item If $M, N \in p$ and $M \cap \omega_1 = \delta_M = \delta_N = N \cap \omega_1$, then $\langle M, \in, <_w \rangle \simeq \langle N, \in, <_w \rangle$
\item If $M \in p$ and $\delta \in dom(p)$ such that $\delta_M < \delta$, then $ \exists N \in p(\delta) \ (M \in N)$.

where $p(\alpha)=\{ M \in p : \delta_M=\alpha\}$ and $dom(p)= \{\alpha: p(\alpha) \neq \emptyset\}$.
\end{itemize}
\end{block}

\end{frame}
\begin{frame}
\frametitle{GITIK'S QUESTION $(2017)$}

Suppose $GCH$ holds and $\kappa$ is a regular cardinal. Is there a cardinal and $GCH$ preserving extension of the universe in which there exists a set $A \subseteq \kappa$ of size $\kappa$ such that for all countable set $X \in \mathscr{P}(\kappa) \cap V$, $A \cap X$ and $X\backslash A$ are non-empty?
\pause
\begin{block}{$\kappa= \aleph_0$}
$\mathbb{P}_{\omega} = \{p:\omega \longrightarrow 2 \ | |p| < \aleph_0 \ \}$
\end{block}
\begin{block}{$\kappa= \aleph_1$}
$\mathbb{P}_{\omega_1} = \{p:\omega_1 \longrightarrow 2 \ | |p| < \aleph_0 \ \}$
\end{block}
\end{frame}
\begin{frame}
\frametitle{DENSITY ARGUMENT}
\begin{block}{}
 $D_X= \{p \in \mathbb{P}: \exists \alpha, \beta \in X \cap dom(p) \big(p(\alpha)=1 \land p(\beta)=0 \big) \}$. 
\end{block}
\begin{block}{}
If $p \in G \cap D_X$ and $\alpha, \beta \in X \cap dom(p)$ be such that $p(\alpha)=1$ and $p(\beta)=0 $, then $\alpha \in X \cap A$ and $\beta \in X \setminus A$
\end{block}
\pause
\begin{block}{THE PROBLEM}
$\mathbb{P}_{\omega_2} \simeq Add(\aleph_0, \aleph_2)$ Hence $2^{\aleph_0}> \aleph_1$ and $GCH$ fails
\end{block}
\end{frame}
\begin{frame}
\frametitle{THE FORCING NOTION}
\begin{definition}
A condition $p$ of $\mathbb{P}$ is a pair $\langle \mathcal{M}_p, f_p \rangle$, whenever:
\begin{itemize}
\item[$(i)$] $\mathcal{M}_p \in \mathbb{P}_{\in}^M$;
\item[$(ii)$] $f_p: \omega_2 \longrightarrow 2$ is a finite partial function;
\item[$(iii)$] If $M,N \in \mathcal{M}_p$ with $\delta_M=\delta_N$, then
 \begin{itemize}
\item $\alpha \in (dom(f_p) \cap M) \Rightarrow \varphi_{M,N}(\alpha) \in dom(f_p)$,
\item for each $\alpha$ as above, $f_p(\varphi_{M,N}(\alpha))=f_p(\alpha)$.
\end{itemize}
\end{itemize}
For $p, q \in \mathbb{P}$, we say $p \leq q$ if and only if $\mathcal{M}_q \subseteq \mathcal{M}_p$ and $f_q \subseteq f_p$.
\end{definition}
\begin{lemma}
$\mathbb{P}$ is strongly proper and satisfies the $\aleph_2$-c.c
\end{lemma}
\begin{lemma}
Forcing with $\mathbb{P}$ preserves the $CH$.
\end{lemma}
\end{frame}
\begin{frame}
\begin{center}
\textbf{{\Huge Thank You}}
\end{center}
\end{frame}
\end{document}